\documentclass[tikz,border=10pt]{standalone}
\input{../../../docs/tikz/dag_preamble.tex}

\begin{document}

\begin{tikzpicture}[
  x=1cm,
  y=1cm,
  every node/.append style={outer sep=4pt},
  source_premise/.style={
    dag_conditional,
    draw=blue!70!black,
    fill=blue!10
  }
]

\dagPaperMetadata
  {Quantifying Spatial Under-reporting Disparities in Resident Crowdsourcing}
  {Liu, Bhandaram, Garg}
  {Nature Computational Science}
  {2024}
  {arXiv:2204.08620 v4 / DOI 10.1038/s43588-023-00572-6}

\dagPaperLegendRightOfMetadata{
\node[dag_result, dag_template_legend] (legRes) at (0,0) {Lean-checked\\result};
\node[source_premise, dag_template_legend] (legModel) at (3.35,0) {Source-model\\premise};
}
\daglegend{(legRes)(legModel)}{Legend}

\begin{dagPaperBody}

\node[source_premise] (ReportingModel) at (0,0) {
  \textbf{Incident/reporting model} \\
  stationary homogeneous incident births; \\
  reports while an incident is active
};

\node[source_premise] (ForwardModel) at (6.5,0) {
  \textbf{Calendar-time first reports} \\
  observed incidents are first reports in \\
  the calendar window; pre-first-report survival
};

\node[dag_result] (PMF) at (13,0) {
  \textbf{Lemma 2 + Condition 1} \\
  selected start is rate-free given the \\
  first report and visible prior history
};

\node[source_premise] (PalmEndpoint) at (19.5,0) {
  \textbf{Condition 2: causal endpoint} \\
  one rate-free endpoint policy may use \\
  visible reports, not future report gaps
};

\node[dag_result] (Lemma1) at (0,-3.5) {
  \textbf{Lemma 1} \\
  duration-mixture first-report probability; \\
  calendar-time first reports are Poisson \\
  at the displayed observed rate
};

\node[dag_result] (Lemma2) at (6.5,-3.5) {
  \textbf{Lemma 2} \\
  wait from the selected start to the next report \\
  has the exponential tail at rate $\lambda$
};

\node[dag_result] (Theorem) at (16.25,-3.5) {
  \textbf{Theorem 1 / Appendix Theorem 2} \\
  $L_\lambda=F\,p(M\mid\lambda(e-s))$; \\
  selected-start likelihood with a causal endpoint
};

\node[dag_result] (Prop1) at (0,-7) {
  \textbf{Proposition 1} \\
  calendar-window counts satisfy the rate limit; \\
  distinct reporting rates can have the same \\
  observed first-report process
};

\node[dag_result] (NYC) at (6.5,-7) {
  \textbf{NYC Eq. (33)} \\
  raw minimum endpoint formula; \\
  stopping-time/no-future property
};

\node[dag_result] (MLE) at (13,-7) {
  \textbf{Eq. (3)} \\
  count/exposure quotient and global MLE, \\
  including the zero-count boundary
};

\node[dag_result] (Chicago) at (19.5,-7) {
  \textbf{Chicago Eq. (34)} \\
  raw minimum of 100-day, closed-time, \\
  and retrieval-time endpoints
};

\node[dag_result] (Regression) at (13,-10.5) {
  \textbf{Eqs. (5)--(7)} \\
  regression log link, proportional likelihood, \\
  and both zero-inflated count branches
};

\draw[dag_arrow] (ReportingModel) -- (Lemma1);
\draw[dag_arrow] (ForwardModel) -- (Lemma1);
\draw[dag_arrow] (Lemma1) -- (Prop1);
\draw[dag_arrow] (PMF) -- (Lemma2);
\draw[dag_arrow] (Lemma2) -- (Theorem);
\draw[dag_arrow] (PalmEndpoint) -- (Theorem);
\draw[dag_arrow] (Theorem) -- (MLE);
\draw[dag_arrow] (MLE) -- (Regression);

\end{dagPaperBody}
\end{tikzpicture}
\end{document}
